\section{Predicates and Quantifiers}

% =========================================================
% TOPIC TITLE
% =========================================================

\begin{frame}{Predicates and Quantifiers}

\begin{center}
\vspace{1cm}
{\Huge\textbf{Predicates and Quantifiers}}

\vspace{0.6cm}


\end{center}

\end{frame}

% =========================================================
% QUESTION 1
% =========================================================

\begin{frame}{Predicates and Quantifiers -- Question 1}

\begin{block}{Question}

A university maintains a database of students enrolled in a
programming course.

Let

$$
P(x):\text{$x$ has submitted the programming assignment}.
$$

Consider the statement:

$$
\text{Every student has submitted the programming assignment.}
$$

\end{block}

\end{frame}

\begin{frame}{Predicates and Quantifiers -- Question 1 (Continued)}

\begin{block}{Question}

\begin{enumerate}
\item Express the statement using a universal quantifier.
\item State what the predicate $P(x)$ represents.
\end{enumerate}

\end{block}

\end{frame}

% ---------------------------------------------------------
% ANSWER 1 -- PART 1
% ---------------------------------------------------------

\begin{frame}{Predicates and Quantifiers -- Question 1: Answer}

\begin{exampleblock}{Step 1: Identify the Predicate}

The predicate is

$$
P(x):\text{$x$ has submitted the programming assignment}.
$$

Here, $x$ represents a student in the course.

\end{exampleblock}

\end{frame}

% ---------------------------------------------------------
% ANSWER 1 -- PART 2
% ---------------------------------------------------------

\begin{frame}{Predicates and Quantifiers -- Question 1: Answer (Continued)}

\begin{exampleblock}{Step 2: Interpret the Predicate}

The predicate becomes either true or false depending on whether a
particular student has submitted the assignment.

The statement says that this property is true for every student.

Therefore, we use the universal quantifier.

\end{exampleblock}

\end{frame}

% ---------------------------------------------------------
% ANSWER 1 -- PART 3
% ---------------------------------------------------------

\begin{frame}{Predicates and Quantifiers -- Question 1: Answer (Conclusion)}

\begin{exampleblock}{Step 3: Apply the Universal Quantifier}

The universal quantifier is

$$
\forall x.
$$

Therefore,

$$
\boxed{\forall x\,P(x)}
$$

where $x$ ranges over all students enrolled in the course.

\end{exampleblock}

\end{frame}

% =========================================================
% QUESTION 2
% =========================================================

\begin{frame}{Predicates and Quantifiers -- Question 2}

\begin{block}{Question}

A cloud service monitors a collection of servers.

Let

$$
P(x):\text{server $x$ is operational}.
$$

The administrator wants to verify whether at least one server is
currently operational.

\end{block}

\end{frame}

\begin{frame}{Predicates and Quantifiers -- Question 2 (Continued)}

\begin{block}{Question}

\begin{enumerate}
\item Express this requirement using an existential quantifier.
\item Explain what the resulting statement means.
\end{enumerate}

\end{block}

\end{frame}

% ---------------------------------------------------------
% ANSWER 2 -- PART 1
% ---------------------------------------------------------

\begin{frame}{Predicates and Quantifiers -- Question 2: Answer}

\begin{exampleblock}{Step 1: Identify the Requirement}

The requirement is that \textbf{at least one} server is operational.

The phrase ``at least one'' indicates the use of the existential
quantifier.

\end{exampleblock}

\end{frame}

% ---------------------------------------------------------
% ANSWER 2 -- PART 2
% ---------------------------------------------------------

\begin{frame}{Predicates and Quantifiers -- Question 2: Answer (Continued)}

\begin{exampleblock}{Step 2: Apply the Existential Quantifier}

The predicate is

$$
P(x):\text{server $x$ is operational}.
$$

The existential quantifier is

$$
\exists x.
$$

Therefore,

$$
\boxed{\exists x\,P(x)}.
$$

\end{exampleblock}

\end{frame}

% ---------------------------------------------------------
% ANSWER 2 -- PART 3
% ---------------------------------------------------------

\begin{frame}{Predicates and Quantifiers -- Question 2: Answer (Conclusion)}

\begin{exampleblock}{Interpretation}

The statement

$$
\exists x\,P(x)
$$

means that there exists at least one server $x$ for which $P(x)$ is
true.

Hence,

$$
\boxed{\text{At least one server is operational.}}
$$

\end{exampleblock}

\end{frame}

% =========================================================
% QUESTION 3
% =========================================================

\begin{frame}{Predicates and Quantifiers -- Question 3}

\begin{block}{Question}

A security system checks every login attempt.

Let

$$
P(x):\text{login attempt $x$ is from an authorized user},
$$

$$
Q(x):\text{login attempt $x$ is successful}.
$$

The security policy states:

$$
\text{Every successful login is from an authorized user.}
$$

\end{block}

\end{frame}

\begin{frame}{Predicates and Quantifiers -- Question 3 (Continued)}

\begin{block}{Question}

Express this policy using predicates, logical implication, and a
universal quantifier.

\end{block}

\end{frame}

% ---------------------------------------------------------
% ANSWER 3 -- PART 1
% ---------------------------------------------------------

\begin{frame}{Predicates and Quantifiers -- Question 3: Answer}

\begin{exampleblock}{Step 1: Translate the Policy}

The statement says that a successful login must be from an authorized
user.

Therefore, for any login attempt $x$,

$$
Q(x)\rightarrow P(x).
$$

\end{exampleblock}

\end{frame}

% ---------------------------------------------------------
% ANSWER 3 -- PART 2
% ---------------------------------------------------------

\begin{frame}{Predicates and Quantifiers -- Question 3: Answer (Continued)}

\begin{exampleblock}{Step 2: Apply the Universal Quantifier}

The policy applies to \textbf{every} login attempt.

Therefore, we use

$$
\forall x.
$$

Combining the quantifier with the implication gives

$$
\boxed{\forall x\,[Q(x)\rightarrow P(x)]}.
$$

\end{exampleblock}

\end{frame}

% ---------------------------------------------------------
% ANSWER 3 -- PART 3
% ---------------------------------------------------------

\begin{frame}{Predicates and Quantifiers -- Question 3: Answer (Conclusion)}

\begin{exampleblock}{Interpretation}

The expression

$$
\forall x\,[Q(x)\rightarrow P(x)]
$$

means that for every login attempt, if the login is successful, then
the user must be authorized.

\end{exampleblock}

\end{frame}

% =========================================================
% QUESTION 4
% =========================================================

\begin{frame}{Predicates and Quantifiers -- Question 4}

\begin{block}{Question}

A machine-learning dataset contains a collection of records.

Let

$$
P(x):\text{record $x$ contains missing values},
$$

$$
Q(x):\text{record $x$ is removed during preprocessing}.
$$

The preprocessing rule is:

$$
\text{Every record containing missing values is removed.}
$$

\end{block}

\end{frame}

\begin{frame}{Predicates and Quantifiers -- Question 4 (Continued)}

\begin{block}{Question}

Express this rule using predicates, implication, and a universal
quantifier.

\end{block}

\end{frame}

% ---------------------------------------------------------
% ANSWER 4 -- PART 1
% ---------------------------------------------------------

\begin{frame}{Predicates and Quantifiers -- Question 4: Answer}

\begin{exampleblock}{Step 1: Identify the Logical Relationship}

The rule states that whenever a record contains missing values, that
record is removed.

Therefore, for any record $x$,

$$
P(x)\rightarrow Q(x).
$$

\end{exampleblock}

\end{frame}

% ---------------------------------------------------------
% ANSWER 4 -- PART 2
% ---------------------------------------------------------

\begin{frame}{Predicates and Quantifiers -- Question 4: Answer (Continued)}

\begin{exampleblock}{Step 2: Apply the Universal Quantifier}

The preprocessing rule applies to every record in the dataset.

Therefore, we use

$$
\forall x.
$$

Combining the quantifier with the implication gives

$$
\boxed{\forall x\,[P(x)\rightarrow Q(x)]}.
$$

\end{exampleblock}

\end{frame}

% ---------------------------------------------------------
% ANSWER 4 -- PART 3
% ---------------------------------------------------------

\begin{frame}{Predicates and Quantifiers -- Question 4: Answer (Conclusion)}

\begin{exampleblock}{Interpretation}

The expression

$$
\forall x\,[P(x)\rightarrow Q(x)]
$$

means that for every record $x$, if $x$ contains missing values, then
$x$ is removed during preprocessing.

\end{exampleblock}

\end{frame}

% =========================================================
% QUESTION 5
% =========================================================

\begin{frame}{Predicates and Quantifiers -- Question 5}

\begin{block}{Question}

A software testing system examines a set of test cases.

Let

$$
P(x):\text{test case $x$ exposes a software defect}.
$$

The testing report states:

$$
\text{No test case exposes a software defect.}
$$

\end{block}

\end{frame}

\begin{frame}{Predicates and Quantifiers -- Question 5 (Continued)}

\begin{block}{Question}

Express this statement using a universal quantifier and predicate
negation.

\end{block}

\end{frame}

% ---------------------------------------------------------
% ANSWER 5 -- PART 1
% ---------------------------------------------------------

\begin{frame}{Predicates and Quantifiers -- Question 5: Answer}

\begin{exampleblock}{Step 1: Interpret the Statement}

The statement says that for every test case, the predicate $P(x)$ is
false.

Therefore, we can express the statement as

$$
\forall x\,\neg P(x).
$$

\end{exampleblock}

\end{frame}

% ---------------------------------------------------------
% ANSWER 5 -- PART 2
% ---------------------------------------------------------

\begin{frame}{Predicates and Quantifiers -- Question 5: Answer (Continued)}

\begin{exampleblock}{Step 2: Interpret the Expression}

The expression

$$
\forall x\,\neg P(x)
$$

means that every test case fails to satisfy the predicate $P(x)$.

In other words, no test case exposes a software defect.

\end{exampleblock}

\end{frame}

% ---------------------------------------------------------
% ANSWER 5 -- PART 3
% ---------------------------------------------------------

\begin{frame}{Predicates and Quantifiers -- Question 5: Answer (Continued)}

\begin{exampleblock}{Step 3: Equivalent Form}

The statement can also be expressed by saying that there does not
exist a test case for which $P(x)$ is true.

Therefore,

$$
\boxed{\forall x\,\neg P(x)
\equiv
\neg\exists x\,P(x)}
$$

\end{exampleblock}

\end{frame}

% ---------------------------------------------------------
% ANSWER 5 -- PART 4
% ---------------------------------------------------------

\begin{frame}{Predicates and Quantifiers -- Question 5: Answer (Conclusion)}

\begin{exampleblock}{Final Interpretation}

Both expressions represent the same statement:

$$
\forall x\,\neg P(x)
$$

and

$$
\neg\exists x\,P(x).
$$

Hence,

$$
\boxed{\text{No test case exposes a software defect.}}
$$

\end{exampleblock}

\end{frame}
